                       IMC 2017, Blagoevgrad, Bulgaria


                                     Day 1, August 2, 2017




Problem 1.    Determine all complex numbers λ for which there exist a positive integer n and
a real n × n matrix A such that A2 = AT and λ is an eigenvalue of A.
                                                                                  (10 points)


Problem 2.    Let f : R → (0, ∞) be a dierentiable function, and suppose that there exists
a constant L > 0 such that
                                         f 0 (x) − f 0 (y) ≤ L x − y
for all x, y . Prove that                           2
                                             f 0 (x) < 2Lf (x)
holds for all x.
                                                                                          (10 points)


Problem 3. For any positive integer m, denote by P (m) the product of positive divisors of m
(e.g. P (6) = 36). For every positive integer n dene the sequence
                     a1 (n) = n,        ak+1 (n) = P (ak (n)) (k = 1, 2, . . . , 2016).

   Determine whether for every set S ⊆ {1, 2, . . . , 2017}, there exists a positive integer n such
that the following condition is satised:
      For every    k with 1 ≤ k ≤ 2017, the number ak (n) is a perfect square if and only if k ∈ S .
                                                                                          (10 points)


Problem 4. There are n people in a city, and each of them has exactly 1000 friends (friendship
is always symmetric). Prove that it is possible to select a group S of people such that at least
n/2017 persons in S have exactly two friends in S .
                                                                                    (10 points)


Problem 5.      Let k and n be positive integers with n ≥ k2 − 3k + 4, and let
                                   f (z) = z n−1 + cn−2 z n−2 + . . . + c0

be a polynomial with complex coecients such that
                                   c0 cn−2 = c1 cn−3 = . . . = cn−2 c0 = 0.

Prove that f (z) and z n − 1 have at most n − k common roots.
                                                                                          (10 points)
